\section{Hierarchy access and complete lifecycle cost}\label{sec:C14:hierarchy}
A hierarchy can keep a small active state, a medium-sized abstraction layer, and a large evidence store. The hierarchy is useful when access frequencies and representation costs differ, but it introduces promotion, eviction, and synchronization policies. If a fact changes, every accessible derivative may need invalidation. If a summary is promoted to active state, it can displace another record whose future value was underestimated. The store layout and the control policy should be analyzed together.

Let $N_r,N_e,N_s$ be the numbers of raw, episodic, and semantic records. A naive dense retrieval over all embeddings has arithmetic cost proportional to $(N_r+N_e+N_s)d$ for embedding dimension $d$, before ranking and model generation. Approximate indexes change the cost-quality tradeoff and add index-maintenance work. A graph hierarchy can reduce some searches while increasing construction and traversal overhead. No storage representation is automatically more efficient without a specified query distribution and update rate.

RecMem's \path{search_memories} exposes three retrieval branches with separately configured limits and thresholds, returning raw, episodic, and semantic results~\cite{repo_recmem2026}. This matters for causal comparison: better answering can arise from more total evidence, different representation, or a different allocation across layers. A baseline that receives fewer retrieved tokens is not a clean test of consolidation structure. Match the complete evidence budget, and separately ablate which layer provides the decisive item.

Construction cost can dominate when many memories are written but few queries arrive. Query-time cost can dominate when a stable store serves many decisions. Let $C_w$ be total writing and consolidation cost, $C_q$ average retrieval-and-generation cost, and $Q\in\mathbb N_{>0}$ the number of useful queries. A simple amortized cost is $(C_w+QC_q)/Q$; it is undefined at $Q=0$, when construction produced no served query. This arithmetic explains why reporting only retrieval tokens can reverse the practical ranking of memory designs. Storage, indexing, failed extractions, and maintenance should remain inside the chosen boundary.

Cache and semantic invalidation also differ. Evicting a vector-search cache can cause recomputation without changing knowledge. Invalidating a false semantic proposition changes what future decisions should use. A hierarchical system must not use the same ``eviction'' operation for both without preserving the semantic reason. The working view should be able to tell whether an item is unavailable temporarily, deleted by retention, superseded by correction, or rejected as untrusted.

The practical implementation should record lifecycle counters by operation: admitted raw records, attempted extractions, successful episodes, semantic revisions, retrieval candidates, final injected items, invalidations, and deletions. These counters are diagnostic observables, not substitutes for task outcomes. They make it possible to explain why a method's cost or reliability changes with the number of sessions rather than reporting a single benchmark aggregate.

Hy-Memory-30B-A3B, introduced in an October 2026 technical report, instead targets a reusable writer inside several existing memory frameworks. It separates writer learning from the fixed reader used for reward, trains on policy-produced bank states, and combines per-chunk probes, attributed cross-chunk feedback, and a bank-growth penalty~\cite{fei2026hymemory}. Its reported cross-harness comparison uses fixed 20\% benchmark subsamples, including a held-out framework. That is evidence about a writer-reader-harness configuration, not a universal ranking of complete memory products. Its learned write protocol also differs from RecMem's training-free recurrence gate.

The comparison identifies a useful causal axis. RecMem changes when expensive abstraction is performed; Hy-Memory changes the policy performing framework-specific writes; HyperSkill changes the representation and maintenance of procedural experience. A result can improve because of any of these interventions, and combining their names does not establish a compatible composite system. A matched study would hold the reader and retrieval configuration fixed while varying admission timing, writer policy, and representation separately, then test interactions. In particular, a writer trained against one retriever can learn representations whose value disappears when the retrieval interface changes.

